\section{An initial pair at the working deletion rate}\label{sec:initial}

This section supplies the induction with a pair whose squared score decays as \(N_0^{-2.20}\) and whose repeated traces have enough common mass to hide later offsets. The initial words have one defect in each alternating period. A finite signed pattern moves some of these defects by one slot in opposite directions. We first describe an observation that forgets the original slot labels, then use its conditional rank probabilities for two purposes: obtaining common mass for pure runs and bounding the score of the signed pair.

\begin{proposition}[Initial pair]\label{prop:initial}
For the parameters of Section~\ref{sec:scales}, there is a choice of the common constants in Definition~\ref{def:invariant} with \(M_B=1\) for which a level-zero pair exists with simulation error \(E_0\le\Delta\). One may choose
\(\omega=\exp(-C_9\Lambda)\) for an absolute constant \(C_9\), and \(C>2\) with \(\log C=O(\Lambda)\).
\end{proposition}

\subsection{Words and the reported mark lists}

In this section \(C_*,c_*>0\) denote absolute constants whose values may change between estimates; constants indexed by \(r\) or \(s_0\) may depend on those fixed indices.

Put \(N=N_0\), \(\eta=.004\), and
\[
 m=\lfloor N^{1/2-\eta}\rfloor,\qquad
 K=\lfloor N^{1/2-2\eta}\rfloor,\qquad
 k=\lfloor\delta m\rfloor,
\]
\[
 J=\left\lfloor
       \frac{\lfloor N/(2m-1)\rfloor-K}{2k+1}
    \right\rfloor,\qquad
 P=J(2k+1)+K.
\]
Lemma~\ref{lem:scale-margins} gives \(\Lambda=o(\log N)\), hence
\(\delta^{-1}=N^{o(1)}\). In particular \(\delta^3m\to\infty\),
\(m\delta^2\to\infty\), and \(K,k+1\le m\) for all sufficiently large instances. The definitions then give
\[
 J\asymp\delta^{-1}N^{2\eta},\qquad J=o(K),\qquad
 P(2m-1)\in[N/2,N].
\]
For the last assertion, \(P\) differs from \(\lfloor N/(2m-1)\rfloor\) by fewer than \(2k+1\), so the omitted length is \(O(km+m)=o(N)\).

For this construction only, let the ordinary block be
\(\mathsf A=01\) and the defect block be \(\mathsf D=1\). For
\(b\in\{-1,0,1\}\), define the period
\[
 W_b=\mathsf A^{\lfloor m/2\rfloor+b-1}\,
       \mathsf D\,
       \mathsf A^{m-\lfloor m/2\rfloor-b}.
\]
It has \(m\) blocks, with the defect at block position
\(\lfloor m/2\rfloor+b\), and length \(2m-1\). Period, slot, and rank indices in this section start at \(1\). For even \(m\), the adjacent types \(W_0,W_1\) are exactly the shifted-defect pair displayed by \citet[arXiv v2, Introduction and Theorem~1]{Chase2021Lower}; for odd \(m\), they are the corresponding smaller pair with one common final \(\mathsf A\) appended. Below we compare \(W_{-1}\) with \(W_1\) in exceptional periods and make signed choices across many periods.

Set
\[
\begin{gathered}
 a=\lfloor P/2\rfloor,\qquad s_0=300,\qquad I=\{0,\ldots,2^{s_0}-1\},\\
 \sum_{i\in I}\beta_i z^i=\prod_{v=0}^{s_0-1}(1-z^{2^v}).
\end{gathered}
\]
The two words \(x_0,y_0\) have \(P\) periods. Every period is \(W_0\) except those at \(a+i\), \(i\in I\), where \(x_0\) uses \(W_{\beta_i}\) and \(y_0\) uses \(W_{-\beta_i}\). For all sufficiently large instances these indices lie in \(\{1,\ldots,P\}\). Both words have length
\(\ell_0=P(2m-1)\in[N/2,N]\), and they differ: the parse into \(01\) and \(1\) from the left is unique. The signs are the classical Prouhet--Thue--Morse signs; their polynomial has a zero of order \(s_0\) at \(1\) \citep[Proposition~2 and Theorem~6]{AlloucheShallit1999}.

We now define the more informative observation used only at this level. For \(H\in\{\mathsf A,\mathsf D\}\), put
\[
 r_e=\delta^2/2,\qquad
 \nu_H=\frac{\Tr(H)-r_e\delta_\varnothing}{1-r_e},
\]
where \(\varnothing\) denotes the empty string. This reserves the same empty output mass in both block laws. In the proof of Lemma~2.1, steps (I) and (III), \citet{HoldenLyons2020} first erase a mass \(q^2/2\) from \(01\) blocks at their deletion rate \(q\); here the same decomposition is applied at rate \(\delta\) to both \(01\) and \(1\), followed by an explicit likelihood score.

For each block independently, silently erase it with probability \(r_e\); otherwise report a mark with law \(\nu_H\). The mark alphabet is
\[
 \mathcal M_{\rm mark}=\{\varnothing,0,1,01\}.
\]
An observation on a sequence of blocks is the ordered list of its returning marks, an element of \(\bigsqcup_{l\ge0}\mathcal M_{\rm mark}^l\). The list reports its individual entries, including empty-valued entries, but no original slot labels or period boundaries. Thus the empty list and the one-entry list \((\varnothing)\) are distinct. Its count \(l\) is the number of entries, not the number of bits in their concatenation. Concatenating the mark strings gives the ordinary trace law, since
\(\Tr(H)=r_e\delta_\varnothing+(1-r_e)\nu_H\).

Under \(\nu_{\mathsf A}\), define
\[
 Z(Y)=\frac{\nu_{\mathsf D}(Y)}{\nu_{\mathsf A}(Y)}-1.
\]
The defect law is dominated by the ordinary-block law. Its support is contained in \(\{\varnothing,1\}\), where the likelihood ratios are
\[
 \frac{\delta-\delta^2/2}{\delta^2/2}
 \quad\text{and}\quad \frac1\delta,
\]
respectively. Hence \(Z\) is centered, \(\|Z\|_\infty=O(\delta^{-1})\), and
\(\|Z\|_{L^2(\nu_{\mathsf A})}^2\lesssim\delta^{-1}\), uniformly for \(\delta\le1/2\).

\subsection{Conditional densities and the rank kernel}

Consider \(1\le u\le m\) periods with fixed types
\(\mathbf b=(b_1,\ldots,b_u)\in\{-1,0,1\}^u\). There are \(n_b=um\) slots, and the defect in period \(j\) is at
\[
 p_{\mathbf b}(j)=(j-1)m+\lfloor m/2\rfloor+b_j.
\]
The return count has law \(\operatorname{Bin}(n_b,1-r_e)\), independent of \(\mathbf b\). Conditional on any count \(0\le l\le n_b\), the returning subset \(S_l\) is uniform among the size-\(l\) subsets of \(\{1,\ldots,n_b\}\). This subset is not reported. If \(p\in S_l\), let \(\operatorname{rk}_{S_l}(p)\) be its rank in the increasing order on \(S_l\). Put \(Q_l=\nu_{\mathsf A}^{\otimes l}\), with \(Q_0\) the point mass on the empty list.

Relative to \(Q_l\), the density of the reported list \(Y=(Y_1,\ldots,Y_l)\) is
\begin{equation}\label{eq:base-density}
 G_{\mathbf b}^{(u,l)}(Y)
 =\mathbb E_{S_l}
   \prod_{\{j:p_{\mathbf b}(j)\in S_l\}}
       \bigl(1+Z(Y_{\operatorname{rk}_{S_l}(p_{\mathbf b}(j))})\bigr).
\end{equation}
Empty products are \(1\), so this also defines the density at \(l=0\). It integrates to one because each \(1+Z\) has mean one under \(\nu_{\mathsf A}\). To expand it, for \(1\le r\le u\) take an increasing tagged tuple
\(\mathbf j=(j_1,\ldots,j_r)\) of periods and an increasing rank tuple
\(1\le v_1<\cdots<v_r\le l\). Write \(p_h=p_{\mathbf b}(j_h)\), and put
\begin{equation}\label{eq:base-rank-event}
\begin{split}
 W_{\mathbf j,\mathbf b}(v)
 &=\Prob\bigl\{
 p_h\in S_l,\ 
 \operatorname{rk}_{S_l}(p_h)=v_h
 \text{ for every }h
 \bigr\},\\
 A_{r,\mathbf b}(v)&=\sum_{1\le j_1<\cdots<j_r\le u}
                         W_{\mathbf j,\mathbf b}(v).
\end{split}
\end{equation}
The event in \(W_{\mathbf j,\mathbf b}\) imposes no restriction on untagged slots, including other defects. It is the joint probability of the tagged returns and ranks within the count conditioning, rather than a rank law conditioned again on all tags returning.
Expanding the product in \eqref{eq:base-density} now gives
\begin{equation}\label{eq:base-density-expansion}
 G_{\mathbf b}^{(u,l)}(Y)
 =1+\sum_{r=1}^{u}\ \sum_{1\le v_1<\cdots<v_r\le l}
       A_{r,\mathbf b}(v)\prod_{h=1}^r Z(Y_{v_h}).
\end{equation}

Define the slot gaps and rank gaps by
\[
\begin{aligned}
 (g_0,\ldots,g_r)
   &=(p_1-1,p_2-p_1-1,\ldots,p_r-p_{r-1}-1,n_b-p_r),\\
 (z_0,\ldots,z_r)
   &=(v_1-1,v_2-v_1-1,\ldots,v_r-v_{r-1}-1,l-v_r).
\end{aligned}
\]
Their sums are \(n_b-r\) and \(l-r\). Put \(\theta=l/n_b\) and
\(f_g(z)=\Prob\{\operatorname{Bin}(g,\theta)=z\}\). Counting returning untagged slots in each gap gives the exact kernel
\begin{equation}\label{eq:base-rank-kernel}
\begin{aligned}
 W_{\mathbf j,\mathbf b}(v)
 &=\frac{\prod_{i=0}^{r}\binom{g_i}{z_i}}{\binom{n_b}{l}}\\
 &=\frac{\theta^r}{f_{n_b}(l)}
       \prod_{i=0}^{r}f_{g_i}(z_i).
\end{aligned}
\end{equation}
We set \(f_g(z)=0\) outside \(z\in\{0,\ldots,g\}\). The last expression represents the same uniform subset by independent Bernoulli-\(\theta\) sampling conditioned on its total, including the degenerate laws at \(\theta=0,1\); it does not change the actual erasure rate \(r_e\).

For the rank estimates, call a count admissible when
\[
 \theta=\frac{l}{n_b}\in[1-\delta^2,1-\delta^2/4].
\]

For a tagged tuple define the coarse gaps
\[
 \gamma_0=j_1-\tfrac12,\qquad
 \gamma_i=j_{i+1}-j_i\ (1\le i<r),\qquad
 \gamma_r=u-j_r+\tfrac12.
\]
They sum to \(u\), and \(g_i=m\gamma_i+\alpha_i\) with
\(|\alpha_i|\le3\). The end coarse gaps are positive half-integers and the internal coarse gaps are positive integers. In particular \(g_i\gtrsim m\); at an admissible count, their expected binomial success and failure counts tend to infinity uniformly by \(m\delta^2\to\infty\).

\begin{lemma}[Rank-kernel square sum]\label{lem:base-kernel}
For the \(u\)-period construction and admissible count above, and every \(1\le r\le u\),
\begin{equation}\label{eq:abase}
 \|A_{r,\mathbf b}\|_{\ell^2}^2
 \le (C_3/\delta)^r(u/m)^{r/2}r^{-r/2},
\end{equation}
where \(\ell^2\) uses counting measure on increasing rank tuples and \(C_3\) is absolute, uniformly in the deterministic period types and the admissible count.
\end{lemma}

\begin{proof}
Fix \(u,l,\mathbf b\); all kernels in this proof therefore share \(\theta\) and \(f_{n_b}(l)\). In an inner product, both kernels use the same rank tuple \(v\), hence the same rank-gap vector \(z(v)\). The correspondence between increasing \(v\) and nonnegative \(z\) with sum \(l-r\) gives
\[
 \langle W_{\mathbf j,\mathbf b},W_{\mathbf j',\mathbf b}\rangle
 =\frac{\theta^{2r}}{f_{n_b}(l)^2}
   \sum_{\substack{z_i\ge0\\\sum_i z_i=l-r}}
       \prod_{i=0}^{r}f_{g_i}(z_i)f_{g_i'}(z_i).
\]
Choose a largest \(\gamma_e\) in the first tuple, breaking ties by the least index. Then \(g_e\gtrsim mu/(r+1)\). Solve the total constraint for \(z_e\), bound that factor by its supremum, and drop the constraints on the other coordinates:
\begin{equation}\label{eq:base-inner-product}
 \langle W_{\mathbf j,\mathbf b},W_{\mathbf j',\mathbf b}\rangle
 \le\frac{\theta^{2r}}{f_{n_b}(l)^2}
      \sup_z f_{g_e}(z)f_{g_e'}(z)
      \prod_{i\ne e}\sum_z f_{g_i}(z)f_{g_i'}(z).
\end{equation}
No multiplicity is introduced when \(z_e\) is solved.

Stirling's formula at the mean \(l=\theta n_b\) gives
\(f_{n_b}(l)^{-2}\lesssim\theta(1-\theta)n_b\lesssim\delta^2n_b\).
Also
\[
 \sup_z f_{g_e}(z)f_{g_e'}(z)\lesssim(\delta^2g_e)^{-1}.
\]
When \(g_e'\ge g_e/2\), this follows from the binomial maximum
\(O((\theta(1-\theta)g)^{-1/2})\) for each factor. When
\(g_e'<g_e/2\), the means differ by a fixed multiple of \(g_e\), so at every common argument one factor is a Hoeffding tail of size
\(\exp(-c_* g_e)\) \cite[Theorem~2]{Hoeffding1963}, which is smaller than the asserted bound. The product of this supremum and the inverse squared denominator is therefore \(O(r+1)\). The Stirling constants are absolute because every relevant expected success and failure count is at least a fixed multiple of \(m\delta^2\to\infty\).

For each free coordinate we use
\begin{equation}\label{eq:base-binomial-convolution}
 \sum_z f_g(z)f_{g'}(z)
 \le (C_{\rm rk}/\delta)g^{-1/2}
       \exp\!\left(-c_{\rm rk}\frac{(g-g')^2}{g+g'}\right)
\end{equation}
with absolute positive constants. To see this, regard the sum as the probability that the difference of two independent binomials is zero. Tilt their product law by \(e^{t\,\mathrm{difference}}\), with
\(t=-e_0(g-g')/(g+g')\) and a sufficiently small absolute \(e_0>0\).
The log moment is
\[
 \theta t(g-g')+O(t^2(g+g'))
 \le-c_{\rm rk}(g-g')^2/(g+g'),
\]
since \(\theta\ge1/2\). Under the tilt independence is preserved, the Bernoulli odds are multiplied by \(e^{\pm t}\), and the failure probabilities remain comparable to \(\delta^2\). The probability of zero under this tilted law is at most the first binomial's maximum mass \(O((\delta\sqrt g)^{-1})\), proving \eqref{eq:base-binomial-convolution}.

Corresponding coarse-gap coordinates have integer differences. Since
\(g_i=m\gamma_i+O(1)\) with an absolute error, for a nonzero mismatch and all large \(m\),
\[
 |g_i-g_i'|\ge\tfrac12m|\gamma_i-\gamma_i'|,
 \qquad g_i+g_i'\lesssim mu,\qquad g_i\gtrsim m\gamma_i.
\]
For a matching coordinate the exponential below is \(1\). Substituting these comparisons into \eqref{eq:base-inner-product} and using \(\theta^{2r}\le1\) yields the intermediate bound
\begin{equation}\label{eq:base-coarse-inner}
 \langle W_{\mathbf j,\mathbf b},W_{\mathbf j',\mathbf b}\rangle
 \le C_{\rm rk}(r+1)
 \prod_{i\ne e}\left[
   \frac{C_{\rm rk}}{\delta\sqrt{m\gamma_i}}\,
   \exp\!\left(-c_{\rm rk}
        \frac{m(\gamma_i-\gamma_i')^2}{u}\right)\right].
\end{equation}

It remains to sum this expression over both tagged tuples. Once the other coarse gaps are fixed, the constraint \(\sum_i\gamma_i=u\) determines the eliminated gap; the same holds for \(\gamma'\). For fixed \(\gamma\), dropping validity of the determined \(\gamma_e'\) can only enlarge the sum, and
\[
 \sum_{\gamma'}\prod_{i\ne e}
       e^{-c_{\rm rk}m(\gamma_i-\gamma_i')^2/u}
 \le\prod_{i\ne e}\sum_{z\in\mathbb Z}
       e^{-c_{\rm rk}(m/u)z^2}
 \le C_{\rm rk}^r
\]
because \(m\ge u\). Thus the remaining \(r\) first-tuple coordinates contribute \(m^{-r/2}\) times a sum of
\(\prod_{i\ne e}\gamma_i^{-1/2}\) with their total at most \(u\).
With \(\tau=r/u\in(0,1]\),
\[
 \sum_{\sum_{i\ne e}\gamma_i\le u}
       \prod_{i\ne e}\gamma_i^{-1/2}
 \le e^{\tau u}\prod_{i\ne e}
       \sum_{\gamma_i>0}\gamma_i^{-1/2}e^{-\tau\gamma_i}
 \le C_{\rm rk}^r(u/r)^{r/2}.
\]
The last lattice sum is \(O(\tau^{-1/2})\) on both the positive integer and positive half-integer lattices, by comparison with the integral of \(x^{-1/2}e^{-\tau x/2}\). Finally sum over the \(r+1\) possible selected coordinates \(e\). The resulting \(O((r+1)^2)\) factor is absorbed into a constant to the power \(r\). Since
\(\|A_{r,\mathbf b}\|_2^2\) is the sum of these inner products, this proves \eqref{eq:abase}.
\end{proof}

\subsection{Common mass for consecutive pure runs}

The rank estimate first gives common mass at the working deletion rate. This mass will later conceal the side orientations that move the exceptional periods as a group.

For \(u=k,k+1\), write \(G^{(u,l)}=G_{\mathbf0}^{(u,l)}\) for the conditional density of the list from \(W_0^u\). At a count \(l\) admissible for both totals, the two densities use the same \(Q_l\), although their respective parameters are \(\theta=l/(um)\). Orthogonality of distinct multilinear monomials in the independent centered scores gives
\[
 \mathbb E_{Q_l}\bigl[(G^{(u,l)})^2\bigr]
 =1+\sum_{r=1}^{u}\|A_{r,\mathbf0}\|_2^2\|Z\|_2^{2r}
 \le1+\sum_{r\ge1}
       (C_6\delta^{-2}\sqrt{u/m})^r r^{-r/2}
 \le\exp(C_6'\Lambda).
\]
Indeed, for \(a\ge0\),
\(\sum_{r\ge1}(a/\sqrt r)^r\le\exp(C_*(1+a^2))\): for \(r>4a^2\) the terms are at most \(2^{-r}\); for \(1\le r\le4a^2\), the maximum of \(r\log(a/\sqrt r)\) over positive reals is at most \(a^2/(2e)\). Here \(\delta^{-4}u/m\lesssim\delta^{-3}\) at both side counts.

Nonnegativity, mean one, and Cauchy--Schwarz imply
\[
 Q_l\{G^{(u,l)}\ge1/2\}\ge\tfrac14\exp(-C_6'\Lambda).
\]
Both densities are nondecreasing functions of the same coordinate likelihood values \(1+Z\ge0\): \eqref{eq:base-density} is an average of nonnegative products in those variables. The coordinates are independent under \(Q_l\). Increasing events for independent real coordinates are nonnegatively correlated, the independent-coordinate case of association \citep[Theorem~2.1]{EPW1967}. For completeness, in one dimension the covariance of increasing bounded functions \(f,g\) equals
\(\mathbb E[(f(X)-f(X'))(g(X)-g(X'))]/2\ge0\) for independent copies \(X,X'\). Induction conditions on the last coordinate and applies the same fact to the conditional means. Applying this to the two threshold events gives
\[
 \int\min(G^{(k,l)},G^{(k+1,l)})\,dQ_l
 \ge\tfrac1{32}\exp(-2C_6'\Lambda).
\]

The count laws overlap on
\[
 |l-(1-r_e)km|\le m.
\]
For either \(n_b=km\) or \(n_b=(k+1)m\), every such count is admissible for all large instances: its distance from the corresponding mean is at most \(2m\), and
\(m/n_b\asymp1/(\delta m)=o(\delta^2)\).
Stirling's formula gives mass per count at least
\[
 c_7(\delta\sqrt{n_b})^{-1}
       \exp\!\left(-C_7\frac{m^2}{\delta^2n_b}\right)
 \ge m^{-1}\exp(-C_8\Lambda)
\]
for absolute positive constants. One way to see the exponent is that the relative entropy between
\(\operatorname{Ber}(l/n_b)\) and \(\operatorname{Ber}(1-r_e)\) is at most a constant times
\((l/n_b-(1-r_e))^2/\delta^2\), while
\[
 (l/n_b)(1-l/n_b)n_b\asymp\delta^2n_b
 \asymp\delta^3m^2\longrightarrow\infty.
\]
There are at least \(m\) integer counts in this interval. Sum the conditional overlap bounds with the lesser of the two count probabilities. For a sufficiently large absolute \(C_9\), the mark-list laws of \(W_0^k,W_0^{k+1}\) have overlap at least
\[
 \omega=\exp(-C_9\Lambda)<1/8.
\]
Their ordinary trace laws have at least the same overlap by concatenation. Choose
\(C=\max\{C_{\rm abs}/\omega,9/\delta\}\). Then \(C\delta>8\) and \(\log C=O(\Lambda)\).

\subsection{The preliminary observation and its high-degree terms}

We use the common mass just obtained for the mark-list laws. Draw \(J\) independent fair bits \(E_1,\ldots,E_J\). Pair a run of \(k+E_i\) periods before a central run of \(K\) periods with a run of \(k+1-E_i\) periods after it; use a fixed order for the runs on each side. The number of periods before the central run is
\[
 L_{\rm left}=Jk+\sum_{i=1}^J E_i.
\]
For every schedule the exceptional global period \(a+i\) has central index \(X+i\), where
\begin{equation}\label{eq:base-anchor}
 X=a-L_{\rm left}
   =\left\lfloor\frac{J+K}{2}\right\rfloor-\sum_{i=1}^J E_i
   =K/2+O(J+1).
\end{equation}
Since \(J=o(K)\) and \(|I|\) is fixed, all exceptional periods lie in the central run for every schedule. Thus each schedule cuts the same prescribed word into pure side runs and a central run with the indicated translated types.

Use the paired flag construction of Section~\ref{sec:tools} on each pair of side mark-list laws, with a common sublaw \(\omega U\) on the list space. Report the flag and both lists; report the orientation bit only on the visible flag. Report one ordered list of individual marks from the central run, without period boundaries or original slot labels. The list values, including empty-valued entries, are retained. The full report also records the fixed left/right order of its side lists.

Let \(P_+,P_-\) be these preliminary report laws for \(x_0,y_0\), and let \(P_{\rm pure}\) be the corresponding law for \(W_0^P\). There is one deterministic projection \(\mathcal K_0\): discard the flags and bits and concatenate all mark values in chronological order. Conditional on any orientation schedule, averaging the pair flags gives independent correct side-list laws; the central list is generated independently from the actual central types. The schedule partitions the same word as just observed. Consequently
\begin{equation}\label{eq:base-exact-projection}
 \mathcal K_0P_+=\Tr(x_0),\qquad
 \mathcal K_0P_-=\Tr(y_0),\qquad
 \mathcal K_0P_{\rm pure}=\Tr(W_0^P).
\end{equation}

For the score analysis, the metadata consist of the side flags and lists, the visible bits, and the central list count \(l\); they exclude the central mark values. Their law \(Q_{\rm meta}\) is common to all three words. The count is independent of the side data and of \(X\), because the \(Km\) central erasures have the same probabilities for every type. Given a metadata value \(\mathfrak d\), let \(M\) be the number of hidden pairs. The remaining hidden bits are independent and fair, so the conditional law \(\pi_{\mathfrak d}\) of \(X\) is a constant minus \(\operatorname{Bin}(M,1/2)\).

Extend \(\beta_i\) by zero outside \(I\), and define the central type vectors by
\((\mathbf b_\pm(X))_j=\pm\beta_{j-X}\). Conditional on \(\mathfrak d\), the central densities and their coefficient arrays are
\begin{equation}\label{eq:base-central-mixture}
\begin{aligned}
 G_\pm^{\mathfrak d}
   &=\mathbb E_{X\sim\pi_{\mathfrak d}}G_{\mathbf b_\pm(X)}^{(K,l)},
 &A_{r,\pm}^{\mathfrak d}
   &=\mathbb E_{X\sim\pi_{\mathfrak d}}A_{r,\mathbf b_\pm(X)},\\
 G_{\rm pure}^{\mathfrak d}&=G_{\mathbf0}^{(K,l)},
 &A_{r,\rm pure}^{\mathfrak d}&=A_{r,\mathbf0}.
\end{aligned}
\end{equation}
These densities all integrate to one under \(Q_l\). In particular
\[
 P_\pm(d\mathfrak d,dY)
 =Q_{\rm meta}(d\mathfrak d)\,
       G_\pm^{\mathfrak d}(Y)Q_l(dY),
\]
with \(l\) read from \(\mathfrak d\), and the analogous identity holds for \(P_{\rm pure}\). This is the common conditional reference required by Lemma~\ref{lem:clipping}.

Call metadata typical when
\[
 \frac l{Km}\in[1-\delta^2,1-\delta^2/4],
 \qquad M\ge\omega^2J/2.
\]
The first failure is a multiplicative binomial tail for the central erasures, and the second is one for the independent hidden flags. The scale margins make both negligible at the required final-length budget:
\[
 (d_D+B+1)\log n=N^{o(\epsilon)},\qquad
 \delta^2Km=N^{1-3\eta-o(1)},\qquad
 \omega^2J=N^{2\eta-o(1)}.
\]
Thus both Chernoff exponents dominate
\(\log\Delta^{-1}=100(d_D+B+1)\log n\), and atypical metadata have mass \(o(\Delta)\).

Fix typical metadata for the remainder of the coefficient analysis and suppress the superscript \(\mathfrak d\). Put
\(\lambda=\sqrt{K/m}=N^{-\eta/2+o(1)}\).
For a fixed central type vector, Lemmas~\ref{lem:base-kernel} and~\ref{lem:moment} give
\[
 \left\|\sum_{1\le v_1<\cdots<v_r\le l}
      A_{r,\mathbf b}(v)\prod_{h=1}^r Z(Y_{v_h})\right\|_{r_0}
 \le \left(C_*r_0\delta^{-3/2}\sqrt\lambda\right)^r r^{-r/4}.
\]
Here \(\|Z\|_{r_0}\lesssim\delta^{-1}\), and
\(r_0,\delta^{-1}=N^{o(1)}\). The factor in parentheses is
\(N^{-\eta/4+o(1)}=N^{-.001+o(1)}\). Put \(R_0=3000\), and for each of the three densities write \(G_{\ge R_0}\) for the sum of its homogeneous terms of degrees at least \(R_0\) in \eqref{eq:base-density-expansion} and \eqref{eq:base-central-mixture}. Geometric summation, followed by Minkowski's inequality when mixing over \(X\), gives
\begin{equation}\label{eq:base-high-degrees}
\begin{aligned}
 \|G-1\|_{r_0}&\le N^{-.0009},\\
 \|G_{\ge R_0}\|_{r_0}&\le N^{-2.3}.
\end{aligned}
\end{equation}
The growing \(r_0\) is a moment order; \(R_0\) is a fixed degree cutoff. It remains to estimate the differences of their finitely many lower-degree arrays.

\subsection{Differences of one rank kernel}

For a tagged tuple in the central \(K\) periods, let \(W_{\mathbf j}\) be its type-\(W_0\) kernel. Let \(T_h\) move the \(h\)-th tagged slot one position to the right, while the rank tuple stays fixed, and put
\[
 D_h=(T_h-T_h^{-1})/2,\qquad
 V_h=(T_h+T_h^{-1}-2)/2.
\]
These operators change only the adjacent slot-gap lengths. In particular the total \(Km\), the count \(l\), the parameter \(\theta=l/(Km)\), and the denominator \(f_{Km}(l)\) remain fixed.

\begin{lemma}[Finite differences of a rank kernel]\label{lem:base-differences}
Fix \(1\le r<R_0\), an admissible central count, and a tagged tuple in \(K\) periods with coarse gaps \(\gamma_i\). For disjoint sets
\(\mathcal D,\mathcal V\subseteq\{1,\ldots,r\}\), put
\(v'=|\mathcal D|+2|\mathcal V|\). Then
\begin{equation}\label{eq:abase-diff}
 \left\|\prod_{h\in\mathcal D}D_h
             \prod_{h\in\mathcal V}V_h W_{\mathbf j}\right\|_{\ell^2}^2
 \le \delta^{-O_r(1)}\sqrt{Km}\,m^{-v'}
             \prod_{i=0}^r(m\gamma_i)^{-1/2}.
\end{equation}
The implicit constants depend only on \(r\).
\end{lemma}

\begin{proof}
Let \(S_i\) increase gap \(i\)'s length in the numerator of
\eqref{eq:base-rank-kernel}. Since \(T_h=S_{h-1}S_h^{-1}\),
\[
\begin{aligned}
 D_h&=\tfrac12(T_h-1)(1+T_h^{-1}),&
 V_h&=\tfrac12(T_h-1)^2T_h^{-1},\\
 T_h-1&=(S_{h-1}-1)S_h^{-1}+(S_h^{-1}-1).
\end{aligned}
\]
Use also \(S_i^{-1}-1=-(S_i-1)S_i^{-1}\). Expanding assigns exactly \(v'\) first length differences to gaps, with \(C_r\) terms and shifts of their lengths by \(O_r(1)\). For an assigned term write \(b_i\) for the order on gap \(i\); thus \(\sum_i b_i=v'\).

Write \((\nabla f)_g=f_{g+1}-f_g\) for the first length difference. At any integer length \(g\ge1\), the order-\(b\) difference \((\nabla^b f)_g\) has Fourier transform
\[
 \theta^b(e^{it}-1)^b(1-\theta+\theta e^{it})^g.
\]
On \([-\pi,\pi]\) the last factor has modulus at most
\(\exp(-c_*\delta^2g t^2)\), since \(\theta(1-\theta)\gtrsim\delta^2\).
Parseval and Fourier inversion therefore give, respectively,
\[
 \|(\nabla^b f)_g\|_2^2
 \le C_b\delta^{-(2b+1)}g^{-b-1/2},\qquad
 \|(\nabla^b f)_g\|_\infty^2
 \le C_b\delta^{-(2b+2)}g^{-b-1}.
\]
All shifted lengths are comparable to \(g_i\) for large \(m\), because \(r\) is fixed and \(g_i\gtrsim m\).

Choose a largest coarse gap \(e\), so its shifted length
\(\widetilde g_e\gtrsim_r Km\). In the squared rank sum, solve for \(z_e\) and use the squared supremum there, with squared \(\ell^2\) sums on the other \(r\) gaps. For one assigned term \(\mathcal F\), the displayed bounds and
\(f_{Km}(l)^{-2}\lesssim\delta^2Km\) give
\[
\begin{aligned}
 \|\mathcal F\|_2^2
 &\le C_r\delta^{-(2v'+r)}Km\,
       \widetilde g_e^{-1-b_e}
       \prod_{i\ne e}\widetilde g_i^{-1/2-b_i}\\
 &\le C_r\delta^{-(2v'+r)}
       \sqrt{Km}\,m^{-v'}
       \prod_{i=0}^r(m\gamma_i)^{-1/2}.
\end{aligned}
\]
The second line uses \(\widetilde g_i\gtrsim m\gamma_i\),
\(\widetilde g_i\gtrsim m\), and the extra factor
\(\widetilde g_e^{-1/2}\lesssim_r(Km)^{-1/2}\).
Combining the \(C_r\) assigned terms proves \eqref{eq:abase-diff}.
\end{proof}

\subsection{Signed coefficient arrays and their sums}

Define the seed motif indicator on the integers by
\(\chi_{\rm seed}(i)={\bf1}_{\{i\in I\}}\). The rank event depends on the type vector only through its tagged positions, so for a fixed \(X\),
\[
 W_{\mathbf j,\mathbf b_\pm(X)}
 =\prod_{h=1}^r T_h^{(\mathbf b_\pm(X))_{j_h}}W_{\mathbf j}.
\]
Each shift factor is exactly
\[
 T_h^{(\mathbf b_\pm(X))_{j_h}}
 =1+\chi_{\rm seed}(j_h-X)V_h
       \ \pm\ \beta_{j_h-X}D_h.
\]
For disjoint \(\mathcal D,\mathcal V\subseteq\{1,\ldots,r\}\), write
\[
 c_{\mathbf j}^{\mathcal D,\mathcal V}
 =\mathbb E_X\left[
       \prod_{h\in\mathcal D}\beta_{j_h-X}
       \prod_{h\in\mathcal V}\chi_{\rm seed}(j_h-X)\right],
 \qquad
 \partial_{\mathcal D,\mathcal V}
 =\prod_{h\in\mathcal D}D_h\prod_{h\in\mathcal V}V_h.
\]
The coefficient differences
\(B_r=A_{r,+}-A_{r,-}\) and
\(B_r^{\rm pure}=A_{r,+}-A_{r,\rm pure}\) come from
\eqref{eq:base-central-mixture}. Here and below, \(\sum_{\mathbf j}\) runs over \(1\le j_1<\cdots<j_r\le K\). Expanding the shift factors gives
\begin{equation}\label{eq:base-pattern-expansion}
 B_r
 =2\sum_{\substack{\mathcal D\cap\mathcal V=\varnothing\\
                    |\mathcal D|\ {\rm odd}}}
       \ \sum_{\mathbf j}
       c_{\mathbf j}^{\mathcal D,\mathcal V}
       \partial_{\mathcal D,\mathcal V}W_{\mathbf j}.
\end{equation}
For \(B_r^{\rm pure}\), the same expansion runs over all nonempty patterns, without the oddness restriction or the factor \(2\). These are finite sums of \(C_r\) patterns because \(r<R_0\) is fixed.

\paragraph{From individual tuples to the full array.}
Within the increasing rank tuples \(1\le v_1<\cdots<v_r\le l\), define
\[
 R_{\mathbf j}
 =\{v:|z_i(v)-\theta m\gamma_i|<m/8
           \text{ for }0\le i<r\}.
\]
These regions are disjoint. The first and internal coarse gaps determine the tuple, and any differing coordinate changes its center by at least \(\theta m\ge3m/4\), more than the combined radii \(m/4\). Bounded shifts of a kernel use the same region.

Fix one pattern and abbreviate its multiplier and operator by
\(c_{\mathbf j}\) and \(\partial\). Disjointness gives the weighted identity
\begin{equation}\label{eq:base-localized-array}
 \left\|\sum_{\mathbf j}{\bf1}_{R_{\mathbf j}}
                  c_{\mathbf j}\partial W_{\mathbf j}\right\|_2^2
 =\sum_{\mathbf j}|c_{\mathbf j}|^2
        \|{\bf1}_{R_{\mathbf j}}\partial W_{\mathbf j}\|_2^2
 \le\sum_{\mathbf j}|c_{\mathbf j}|^2
        \|\partial W_{\mathbf j}\|_2^2.
\end{equation}
The part outside these regions is negligible as an analytic coefficient estimate. To estimate it, expand \(D_h,V_h\) directly into tag shifts \(T_h\), which preserve the total \(Km\). Each resulting term is a fixed signed coefficient times a nonnegative shifted rank kernel. For each such kernel in \(\partial W_{\mathbf j}\), the conditional binomial formula and Hoeffding's inequality \cite[Theorem~2]{Hoeffding1963} give
\[
 \sum_{v\notin R_{\mathbf j}}W_{\mathbf j,\mathrm{shift}}(v)
 \le C_r f_{Km}(l)^{-1}e^{-c_* m/K}
 \le C_r\sqrt{Km}\,e^{-c_* m/K}.
\]
Indeed a bounded shift changes each relevant mean by \(O_r(1)\), while the excluded deviation is order \(m\), and each gap has length at most \(Km+O_r(1)\). Dropping the total constraint and taking a union bound over \(i<r\) proves the display. There are at most \(K^r\) tuples, the multiplier is bounded by \(1\), and \(\partial\) is a fixed linear combination of shifted kernels. Using \(\|\cdot\|_2\le\|\cdot\|_1\) for coefficient arrays, the discarded array has \(\ell^2\) norm at most
\[
 \tau_r=C_rK^r\sqrt{Km}\,e^{-c_* m/K}.
\]
It follows from \eqref{eq:base-localized-array} that the squared norm of the full array for this pattern is at most
\[
 2\sum_{\mathbf j}|c_{\mathbf j}|^2
        \|\partial W_{\mathbf j}\|_2^2+2\tau_r^2.
\]
This restriction does not condition a reported law or incur a total-variation error. Since \(m/K=N^{\eta+o(1)}\), even
\((C_*r_0\delta^{-1})^r\tau_r=o(N^{-A})\) for every fixed \(A\), uniformly over the finitely many \(r<R_0\).

\paragraph{Sums with one involved period fixed.}
The next two sums separate rank differentiation from the random location multiplier. For \(1\le h\le r\) and an integer \(j\), put
\[
 \Sigma_{r,h}(j)=\sqrt{Km}
 \sum_{\substack{1\le j_1<\cdots<j_r\le K\\j_h=j}}
             \prod_{i=0}^r(m\gamma_i)^{-1/2}.
\]
For \(r\ge2\), let \(\Sigma_{r,h}^{\rm bd}(j)\) be the same sum restricted to tuples with \(\gamma_i\le|I|\) for some \(1\le i<r\).

\begin{lemma}[Gap sums at a fixed period]\label{lem:base-fixed-sums}
For fixed \(1\le r<R_0\), uniformly in \(h\) and
\(j\in[K/3,2K/3]\),
\begin{equation}\label{eq:base-fixed-sums}
 \Sigma_{r,h}(j)\le C_r(Km)^{-1/2}\lambda^{r-1},\qquad
 \Sigma_{r,h}^{\rm bd}(j)
 \le C_{r,s_0}(Km)^{-1/2}m^{-1/2}\lambda^{r-2}\quad(r\ge2).
\end{equation}
\end{lemma}

\begin{proof}
The left gaps \(0,\ldots,h-1\) sum to \(j-\tfrac12\), and the right gaps \(h,\ldots,r\) sum to \(K-j+\tfrac12\). Each total is order \(K\). Select a largest gap on each side; both are at least \(c_rK\), so their two factors cost \(C_r(Km)^{-1}\). Given \(j\) and the other \(r-1\) gaps, the two total constraints determine the selected gaps. Each remaining free sum costs
\[
 \sum_{\substack{\gamma>0\\\gamma\le K}}(m\gamma)^{-1/2}
 \le C_*\sqrt{K/m}=C_*\lambda
\]
on either the integer or half-integer lattice. Grouping by the at most \(C_r\) choices of selected gaps and multiplying by \(\sqrt{Km}\) proves the first bound.

For the restricted sum, an internal gap at most \(|I|\) cannot be either selected largest gap for all sufficiently large \(K\). Union over its at most \(r-1\) possible indices. Its sum is at most
\(\sum_{\gamma=1}^{|I|}(m\gamma)^{-1/2}\le C_{s_0} m^{-1/2}\)
instead of \(C_*\lambda\). The other \(r-2\) free sums are unchanged, proving the second bound.
\end{proof}

\paragraph{The location multipliers.}
For the fixed metadata, \eqref{eq:base-anchor} shows that every
\(j\in\operatorname{supp}(\chi_{\rm seed}*\pi)\) satisfies
\(j=K/2+O(J+|I|)\), and hence lies in \([K/3,2K/3]\) for all large instances. Thus Lemma~\ref{lem:base-fixed-sums} applies to every involved period. For a sequence \(a\) on the integers, use the Fourier transform
\(\widehat a(t)=\sum_j a(j)e^{ijt}\). If \(M\ge\omega^2J/2\) is the hidden count, then
\[
 |\widehat\pi(t)|=|\cos(t/2)|^M\le e^{-c_*Mt^2}
 \qquad(-\pi\le t\le\pi).
\]
Since
\[
 \left|\sum_{i\in I}\beta_i e^{iti}\right|
 =\left|\prod_{v=0}^{s_0-1}(1-e^{i2^v t})\right|
 \le C_{s_0}|t|^{s_0},
\]
Parseval and the binomial decay give
\begin{equation}\label{eq:base-multiplier-norms}
\begin{aligned}
 \|\beta*\pi\|_2^2
 &\le C_{s_0} M^{-s_0-1/2}
 \le C_{s_0}\omega^{-2s_0-1}J^{-s_0-1/2},\\
 \|\chi_{\rm seed}*\pi\|_2^2
 &\le \|\chi_{\rm seed}\|_1^2\|\pi\|_2^2
 \le C_{s_0} M^{-1/2}
 \le C_{s_0}\omega^{-1}J^{-1/2}.
\end{aligned}
\end{equation}
The constants depending on \(s_0=300\) are fixed.

The dyadic signs turn convolution into a high-order finite difference. The same cancellation algebra appears in
\citet[Section~3.2, equation~(3.10)]{HMPW2008}, where iterated input-position differences act on a mean-coordinate binomial kernel. Here the signs are convolved with the conditional location law inside a likelihood coefficient; the terms involving several scores require the separate gap estimate above.

In \eqref{eq:base-pattern-expansion}, a singleton
\(\mathcal D=\{h\},\mathcal V=\varnothing\) has \(v'=1\) and multiplier
\((\beta*\pi)(j_h)\). Every other surviving pattern has an odd positive number of \(D\)'s, so \(v'\ge3\) and at least two indices are involved. At any chosen involved \(h\),
\[
 |c_{\mathbf j}^{\mathcal D,\mathcal V}|
 \le(\chi_{\rm seed}*\pi)(j_h).
\]
A nonzero multiplier requires all involved periods to lie in a common translate of \(I\). Two of them therefore differ by at most \(|I|-1\), so some intervening internal gap is at most \(|I|\).

For either kind of pattern choose such an involved \(h\), and let
\(\psi(j)\) be the corresponding envelope
\(|(\beta*\pi)(j)|\) or \((\chi_{\rm seed}*\pi)(j)\).
Apply Lemma~\ref{lem:base-differences} and first sum over tuples with \(j_h=j\):
\[
 \sum_{\mathbf j}|c_{\mathbf j}|^2
        \|\partial W_{\mathbf j}\|_2^2
 \le \delta^{-O_r(1)}m^{-v'}
       \sum_j\psi(j)^2
       \begin{cases}
          \Sigma_{r,h}(j),&\text{for a singleton},\\
          \Sigma_{r,h}^{\rm bd}(j),&\text{otherwise}.
       \end{cases}
\]
Now \eqref{eq:base-fixed-sums} and \eqref{eq:base-multiplier-norms} apply. Combining the finitely many patterns in
\eqref{eq:base-pattern-expansion} and the analytic tail gives
\begin{equation}\label{eq:base-low-coefficients}
\begin{aligned}
 \|B_r\|_2^2\le C_{r,s_0}\delta^{-O_r(1)}\Big[
 &\omega^{-2s_0-1}J^{-s_0-1/2}m^{-1}(Km)^{-1/2}\lambda^{r-1}\\
 &+{\bf1}_{r\ge2}\,\omega^{-1}J^{-1/2}
           (Km)^{-1/2}m^{-7/2}\lambda^{r-2}\Big]+C_r\tau_r^2.
\end{aligned}
\end{equation}
The nonlinear \(m^{-7/2}\) is \(m^{-3}\) from its minimum difference order and \(m^{-1/2}\) from the bounded internal gap. For \(B_r^{\rm pure}\), the singleton \(D\) contribution is the same. Every other pattern has \(v'\ge2\), but a singleton \(V\) requires no bounded internal gap. The full pure-comparison array therefore satisfies
\begin{equation}\label{eq:base-pure-coefficients}
\begin{aligned}
 \|B_r^{\rm pure}\|_2^2\le C_{r,s_0}\delta^{-O_r(1)}\Big[
 &\omega^{-2s_0-1}J^{-s_0-1/2}m^{-1}(Km)^{-1/2}\lambda^{r-1}\\
 &+\omega^{-1}J^{-1/2}
       (Km)^{-1/2}m^{-2}\lambda^{r-1}\Big]+C_r\tau_r^2.
\end{aligned}
\end{equation}

\subsection{Closing the initial invariant}

For each fixed \(r<R_0\), Lemma~\ref{lem:moment} converts the coefficient estimate to the required score estimate:
\[
 \left\|\sum_{1\le v_1<\cdots<v_r\le l}
       B_r(v)\prod_{h=1}^r Z(Y_{v_h})\right\|_{r_0}^2
 \le (C_1r_0\|Z\|_{r_0})^{2r}\|B_r\|_2^2.
\]
The analogous inequality applies to the pure comparison. The fixed pattern constants, the displayed powers of \(\delta^{-1}\) and \(\omega^{-1}\), and the moment factors are \(N^{o(1)}\): use
\(\log r_0=O_B(D+1)\), \(\log\omega^{-1}=O(\Lambda)\), and
\(\Lambda=o(\log N)\). The analytic tails remain below every fixed power after these factors.

The worst nonlinear term of \eqref{eq:base-low-coefficients}, at \(r=2\), is
\[
 J^{-1/2}(Km)^{-1/2}m^{-7/2}
 =N^{-\eta}\,N^{-1/2+3\eta/2}\,
       N^{-7/4+7\eta/2+o(1)}
 =N^{-9/4+4\eta+o(1)}=N^{-2.234+o(1)}.
\]
The singleton term is \(N^{-3.394+o(1)}\) at its worst degree \(r=1\). The worst nonlinear term in \eqref{eq:base-pure-coefficients} is
\(N^{-3/2+5\eta/2+o(1)}=N^{-1.49+o(1)}\).
Sum over the finitely many lower degrees and add the degree tails in \eqref{eq:base-high-degrees}. On every typical metadata value this yields
\begin{equation}\label{eq:base-score-differences}
 \|G_+-G_-\|_{r_0}^2\le N^{-2.22},\qquad
 \|G_+-G_{\rm pure}\|_{r_0}^2\le N^{-1.4}.
\end{equation}

Let \(\mathcal A\) be the joint event of typical metadata and
\(G_+^{\mathfrak d}(Y)\ge1/2\), and define
\[
 \mu_0=P_+(\,\cdot\mid\mathcal A),\qquad
 \nu_0=P_-(\,\cdot\mid\mathcal A).
\]
Lemma~\ref{lem:clipping} applies to the common conditional reference above: \eqref{eq:base-high-degrees} gives
\(\rho=.0009\ge2^{-4B-10}\), and \(2.22>2.20+.01\) supplies its difference margin. Hence
\[
 \nu_0\ll\mu_0,\qquad
 \|d\nu_0/d\mu_0-1\|_{L^{r_0}(\mu_0)}^2\le N^{-2.20},
\]
and conditioning changes each preliminary law by \(o(\Delta)\) in total variation. By \eqref{eq:base-exact-projection} and contraction under the same \(\mathcal K_0\),
\[
 \TV(\mathcal K_0\mu_0,\Tr(x_0))\le o(\Delta)\le\Delta,\qquad
 \TV(\mathcal K_0\nu_0,\Tr(y_0))\le o(\Delta)\le\Delta
\]
for all large instances. This gives \(E_0\le\Delta\) and the score part of the invariant. The hidden-bit law was used for the preliminary densities only.

It remains to transfer pure-run overlap to \(x_0\). This comparison uses the preliminary laws, rather than the full clipping conclusion for their weaker \(N^{-1.4}\) difference. On typical metadata and \(G_+\ge1/2\),
\[
 (\sqrt{G_+}-\sqrt{G_{\rm pure}})^2
 \le 2|G_+-G_{\rm pure}|^2.
\]
On the complement use the bound \(G_++G_{\rm pure}\). The high-moment estimate \eqref{eq:base-high-degrees} and the clipping proof bound the mass of the bad-density region under both preliminary laws by \(o(\Delta)\); atypical metadata have the same bound. Thus \eqref{eq:base-score-differences} gives
\[
 H^2(P_+,P_{\rm pure})\le2N^{-1.4}+o(\Delta)\le3N^{-1.4}.
\]
Projection by \(\mathcal K_0\) gives the same upper bound for
\(H^2(\Tr(x_0),\Tr(W_0^P))\).

Set \(k_0=\lceil CkP\rceil\). Since \(k\sim\delta m\) and
\(C\delta>8\), we have \(k_0\ge\ell_0\) for all large instances. Also,
\[
 \frac{k_0}{\ell_0}
 \le\frac{C\delta m}{2m-1}+\frac1{\ell_0}<C,
\]
using \(\delta\le1/2\), \(C\ge18\), and \(\ell_0\to\infty\). Thus the common choice \(M_B=1\) meets the level-zero count predicate. Moreover
\(C=\exp(O(\Lambda))=N^{o(1)}\). For \(F=k_0,k_0+1\), tensorization followed by concatenation therefore gives
\[
 \TV(\Tr(x_0^F),\Tr(W_0^{PF}))
 \le\sqrt{F\,H^2(\Tr(x_0),\Tr(W_0^P))}
 \le\sqrt{3F N^{-1.4}}
 =N^{-0.2+o(1)}=o(1).
\]
Lemma~\ref{lem:lift}, with base word \(W_0\), overlap count \(k\), replication \(P\), and \(F=k_0\), gives overlap at least \(2/3\) for the two pure repeated comparison words. Paying the two replacement costs leaves
\[
 \ov(\Tr(x_0^{k_0}),\Tr(x_0^{k_0+1}))
 \ge2/3-o(1)\ge\omega.
\]
Together with the length, distinction, score, and simulation estimates already proved, this completes Proposition~\ref{prop:initial}.
